\section*{Devoir maison }

On s'intéresse à la fonction verifie qui prend en paramètre un tableau de valeurs numériques non vide et qui renvoie True si ce tableau est trié dans l’ordre croissant, False sinon.

\begin{lstlisting}[language=Python]
def verifie(tab):
    for i in range(len(tab)-1):
        if tab[i] > tab[i+1]:
            return False
        return True
\end{lstlisting}

\begin{Exercise}[title={Vérification tri}]
\Question Combien y a-t-il de structures de contrôle dans le code ci-dessus ?\\

\Question Quel type de parcours est utilisé à la ligne 2 de la fonction ?\\
\Question Compléter le tableau d'avancement des variables pour un appel de la fonction verifie([0, 5, 8, 8, 9]).\\

\begin{center}
\large
\begin{tabular}{c|c|c|c|c|}
\textbf{Tour} & \textbf{i} & \textbf{tab[i]} & \textbf{tab[i+1]} & \textbf{tab[i] > tab[i+1]}  \\
\hline
0 &  &  &  & \\
1 &  &  &  & \\
2 &  &  &  & \\
3 &  &  &  & \\
4 &  &  &  & \\
5 &  &  &  & \\
\end{tabular}
\end{center}

\Question Quelles sont les valeurs de retours de la fonction verifie pour les appels suivants : \\
\begin{itemize}
    \item verifie([8, 12, 4])
    \item verifie([-1, 4])
    \item verifie([5])
\end{itemize}

On souhaite maintenant modifier la fonction en y rajoutant un paramètre ordre permettant de choisir l'ordre dans lequel le tableau doit être trié. Le paramètre ordre prend deux valeurs possibles : "croissant" ou "décroissant". Par défaut, il a la valeur "croissant".\\

\Question Quelles sont les valeurs de retours de la nouvelle fonction verifie pour les appels suivants : \\
\begin{itemize}
 \item verifie([-1, 4, 8, 8, 88])
 \item verifie([5, 3, 1], ordre="decroissant")
 \item verifie([5, 3, 1])
 \item verifie([7, 8, 9], ordre="decroissant")
\end{itemize}

\Question Compléter le code ci-dessous pour obtenir ce nouveau comportement de la fonction verifie.

\begin{lstlisting}[language=Python]
def verifie(tab, ordre="croissant"):
    .............................................................
    .............................................................
    .............................................................
    .............................................................
    .............................................................
    .............................................................
    .............................................................
    .............................................................
    .............................................................
\end{lstlisting}
\end{Exercise}